\section*{Capítulo \ref{ch:b1:enzymes} --- Enzimas e catálise bioquímica}

\begin{solution}{exo:b1:enzymes:1}
Ela baixa a \omterm{prop:b1:enzymes:whatitdoes}{energia de ativação} (a altura do \omterm{prop:b1:enzymes:whatitdoes}{estado de transição}) e, com
isso, a velocidade nos dois sentidos; ela deixa inalterados o
$\Delta G$, a constante de equilíbrio e a posição do equilíbrio. No
perfil, o pico é rebaixado, e as duas pontas não.
\end{solution}

\begin{solution}{exo:b1:enzymes:2}
$K_m$: concentração de \omterm{def:b1:enzymes:enzyme}{substrato} na metade de $V_{\max}$
(\unit{mol/L}). $V_{\max}$: velocidade com \omterm{def:b1:enzymes:enzyme}{substrato} saturante
(\unit{mol/s}, ou \unit{\micro
mol/min}). $k_{\text{cat}}$: moléculas de \omterm{def:b1:enzymes:enzyme}{substrato} convertidas por
segundo e por \omterm{def:b1:enzymes:enzyme}{enzima} na saturação (\unit{s^{-1}}).
$k_{\text{cat}}/K_m$: eficiência, a constante de velocidade de segunda
ordem em \omterm{def:b1:enzymes:enzyme}{substrato} baixo (\unit{L.mol^{-1}.s^{-1}}).
\end{solution}

\begin{solution}{exo:b1:enzymes:3}
Sem \omterm{def:b1:enzymes:inhibitors}{inibidor}: interseção 1, logo $V_{\max} = 1$; interseção em $x$
$-0.5$, logo $K_m = 2$. Competitivo: $V_{\max} = 1$, interseção em $x$
$-0.25$, $K_m = 4$. Não competitivo: interseção 2, $V_{\max} = 0.5$;
interseção em $x$ $-0.5$, $K_m = 2$.
\end{solution}

\begin{solution}{exo:b1:enzymes:4}
Regulação \omterm{def:b1:proteins:allostery}{alostérica} (por retroalimentação), milissegundos; \omterm{prop:b1:enzymes:control}{modificação
covalente} (fosforilação), de segundos a minutos; ativação proteolítica
de um zimogênio, uma única vez, sob demanda; controle da quantidade pela
expressão gênica e pela degradação, de minutos a horas.
\end{solution}

\begin{solution}{exo:b1:enzymes:5}
$v_0 = 50[S]/(0.1 + [S])$: 8.3, 25, 45.5 e \qty{49.5}{\micro mol/min}.
$0.9 = [S]/(0.1 + [S])$ dá $[S] = 0.9\,K_m/0.1 = \qty{0.9}{mmol/L}$.
\end{solution}

\begin{solution}{exo:b1:enzymes:6}
$e^{20\,000/2580} = e^{7.75} = 2300$. Para $10^{10}$: $\Delta E_a =
RT\ln 10^{10} = 2.58\times 23.0 = \qty{59}{kJ/mol}$.
\end{solution}

\begin{solution}{exo:b1:enzymes:7}
$V_{\max} = \qty{0.12}{mol/L}$ por minuto em \qty{1}{mL}:
\qty{1.2e-4}{mol/min} $= \qty{2.0e-6}{mol/s}$ de produto; \omterm{def:b1:enzymes:enzyme}{enzima}
\qty{2e-12}{mol}: $k_{\text{cat}} = 2\times 10^{-6}/2\times 10^{-12} =
\qty{1e6}{s^{-1}}$. Eficiência $10^6/0.012 = \qty{8e7}{L.mol^{-1}.s^{-1}}$,
a menos de um fator dez do limite da difusão: quase todo encontro com
uma molécula de \omterm{def:b1:enzymes:enzyme}{substrato} é produtivo.
\end{solution}

\begin{solution}{exo:b1:enzymes:8}
$1 + 2/K_i = 2$: $K_i = \qty{2}{mmol/L}$. Para \qty{90}{\%}:
$[S] = 9K_m^{\text{ap}}$: \qty{9}{mmol/L} sem o \omterm{def:b1:enzymes:inhibitors}{inibidor},
\qty{18}{mmol/L} com ele.
\end{solution}

\begin{solution}{exo:b1:enzymes:9}
Inibir a primeira etapa comprometida detém a via inteira de uma vez e
não desperdiça intermediários, que de outro modo se acumulariam; o
produto final é o sinal que importa --- a abundância dele é exatamente a
informação de que a via já não é necessária ---, ao passo que o nível de
um intermediário nada diz sobre a demanda.
\end{solution}

\begin{solution}{exo:b1:enzymes:10}
Uma protease ativa na \omterm{def:b1:cell-unit-of-life:cell}{célula} que a fabrica digeriria essa \omterm{def:b1:cell-unit-of-life:cell}{célula}; feita
como zimogênio inativo, ela é inofensiva até que a enteropeptidase, uma
\omterm{def:b1:enzymes:enzyme}{enzima} do revestimento intestinal, corte o tripsinogênio em tripsina no
duodeno, onde a digestão é desejada; a tripsina então ativa os demais
zimogênios (uma cascata). Como um traço de tripsina pode disparar a
cascata prematuramente, o pâncreas empacota nos grânulos dele um
\omterm{def:b1:enzymes:inhibitors}{inibidor} específico de tripsina como seguro; a falha desse arranjo é a
pancreatite.
\end{solution}

\begin{solution}{exo:b1:enzymes:11}
Sem ele: $3^3/(3^3 + 3^3) = 0.50$. Com ele: $27/(166 + 27) = 0.14$: a
velocidade cai a \qty{28}{\%} do valor dela. Uma \omterm{def:b1:enzymes:enzyme}{enzima} de
Michaelis--Menten com o $K_m$ elevado de 3 para 5.5: de $3/6 = 0.50$ a
$3/8.5 = 0.35$, ou seja, \qty{70}{\%}. A \omterm{def:b1:proteins:allostery}{cooperatividade} torna o mesmo
deslocamento da curva mais que duas vezes mais eficaz perto do ponto de
trabalho: um interruptor, e não um dimmer.
\end{solution}

\begin{solution}{exo:b1:enzymes:12}
Um catalisador nu (uma superfície de platina, um ácido) acelera muitas
reações indiscriminadamente; uma \omterm{def:b1:enzymes:enzyme}{enzima} acelera uma reação sobre um
\omterm{def:b1:enzymes:enzyme}{substrato} --- a especificidade do \omterm{def:b1:enzymes:enzyme}{sítio ativo} dela já é informação sobre
o que a \omterm{def:b1:cell-unit-of-life:cell}{célula} quer que seja feito. A regulação acrescenta uma segunda
camada: os sítios alostéricos, os sítios de fosforilação e os peptídeos
dos zimogênios dizem à \omterm{def:b1:enzymes:enzyme}{enzima} quando e onde agir, em resposta ao estado
da \omterm{def:b1:cell-unit-of-life:cell}{célula}. O custo é uma \omterm{def:b1:proteins:peptide}{proteína} grande (centenas de resíduos para um
sítio de uma dúzia), uma renovação mais lenta que a dos melhores
catalisadores inorgânicos, e a maquinaria para fabricá-la, modificá-la e
destruí-la; o benefício é uma química que só funciona onde e quando é
útil, que é o que um metabolismo é.
\end{solution}

\begin{solution}{pb:b1:enzymes:1}
\textbf{1.} $1/[S]$: 2, 1, 0.5, 0.2, 0.1, 0.05. $1/v_0$: 5.0, 3.0, 2.0,
1.4, 1.2, 1.1.
\textbf{2.} Cada passo de $1/[S]$ muda $1/v_0$ na mesma proporção:
inclinação $(5.0 - 1.1)/(2 - 0.05) = 2.0$; interseção $1.0$.
\textbf{3.} $V_{\max} = 1/1.0 = \qty{1.0}{\micro mol/min}$; $K_m =
\text{inclinação}\times V_{\max} = \qty{2.0}{mmol/L}$.
\textbf{4.} $1.0\times 5/(2 + 5) = 0.714$: como medido.
\textbf{5.} $\qty{10e-9}{mol/L}\times\qty{1e-3}{L} = \qty{1e-11}{mol}$;
$V_{\max} = \qty{1e-6}{mol/min} = \qty{1.67e-8}{mol/s}$;
$k_{\text{cat}} = 1.67\times 10^{-8}/10^{-11} = \qty{1670}{s^{-1}}$.
\textbf{6.} $1670/(2\times 10^{-3}) = \qty{8.3e5}{L.mol^{-1}.s^{-1}}$:
mil vezes abaixo do limite da difusão; só uma colisão em mil é
produtiva.
\textbf{7.} $0.95 = [S]/(2 + [S])$: $[S] = 19\,K_m = \qty{38}{mmol/L}$.
\textbf{8.} $1/v_0$: 9.0, 5.0, 3.0, 1.8, 1.4, 1.2. Inclinação $(9.0 -
1.2)/1.95 = 4.0$; interseção $1.0$.
\textbf{9.} $V_{\max}^{\text{ap}} = \qty{1.0}{\micro mol/min}$
(inalterado); $K_m^{\text{ap}} = 4.0\times 1.0 = \qty{4.0}{mmol/L}$.
\textbf{10.} $K_m$ dobrado, $V_{\max}$ inalterado: competitivo.
\textbf{11.} $4 = 2(1 + 1/K_i)$: $K_i = \qty{1.0}{mmol/L}$.
\textbf{12.} Com o \omterm{def:b1:enzymes:inhibitors}{inibidor}, $v = V_{\max}[S]/(K_m(1 + [I]/K_i) +
[S])$; reduzir a velocidade à metade exige $K_m(1 + [I]/K_i) + [S] = 2(K_m +
[S])$, ou seja, $[I] = K_i(K_m + [S])/K_m$: a \qty{2}{mmol/L}, $[I] =
1\times 4/2 = \qty{2}{mmol/L}$; a \qty{20}{mmol/L}, $[I] = 22/2 =
\qty{11}{mmol/L}$. Quanto mais \omterm{def:b1:enzymes:enzyme}{substrato}, mais \omterm{def:b1:enzymes:inhibitors}{inibidor} é preciso.
\textbf{13.} Uma molécula parecida com o \omterm{def:b1:enzymes:enzyme}{substrato} (ou com o \omterm{prop:b1:enzymes:whatitdoes}{estado de
transição} dele) --- um análogo de éster que não pode ser hidrolisado ---
que se liga no \omterm{def:b1:enzymes:enzyme}{sítio ativo}.
\textbf{14.} $1/v_0$: 10.0, 6.0, 4.0, 2.8, 2.4, 2.2; inclinação $4.0$,
interseção $2.0$: $V_{\max}^{\text{ap}} = \qty{0.5}{\micro mol/min}$,
$K_m^{\text{ap}} = 4.0\times 0.5 = \qty{2.0}{mmol/L}$.
\textbf{15.} $V_{\max}$ reduzido à metade, $K_m$ inalterado: não
competitivo; $2 = 1 + 1/K_i$: $K_i = \qty{1.0}{mmol/L}$.
\textbf{16.} J se liga a um sítio diferente do ativo, tanto na \omterm{def:b1:enzymes:enzyme}{enzima}
livre quanto no $ES$, com a mesma afinidade: o \omterm{def:b1:enzymes:enzyme}{substrato} não compete com
ele, e em qualquer $[S]$ metade das moléculas de \omterm{def:b1:enzymes:enzyme}{enzima} fica inativada.
\textbf{17.} Com J a velocidade dobra (metade de uma quantidade dobrada
de \omterm{def:b1:enzymes:enzyme}{enzima} continua ativa); com I ela também dobra (a velocidade é
proporcional à \omterm{def:b1:enzymes:enzyme}{enzima} em todo $[S]$): dobrar a \omterm{def:b1:enzymes:enzyme}{enzima} nunca distingue os
dois \omterm{def:b1:enzymes:inhibitors}{inibidores}.
\textbf{18.} Velocidade a \qty{25}{\celsius}: $V_{\max}[S]/(K_m + [S])$
com $V_{\max} = k_{\text{cat}}[E]$: por litro, $1670\times 10^{-8} =
\qty{1.67e-5}{mol/L/s}$; $\times 0.5/2.5 = \qty{3.3e-6}{mol/L/s}$; a
\qty{37}{\celsius}, $\times 2^{1.2} = 2.3$: \qty{7.7e-6}{mol/L/s}; em
\qty{1e-12}{L}: \qty{7.7e-18}{mol/s} $= \num{4.6e6}$ moléculas por
segundo.
\textbf{19.} Não é preciso elevar: a velocidade (\num{4.6e6}) supera a
demanda (\num{3e6}); se fosse preciso elevá-la, a \omterm{def:b1:cell-unit-of-life:cell}{célula} poderia
aumentar o \omterm{def:b1:enzymes:enzyme}{substrato}, ou fosforilar a \omterm{def:b1:enzymes:enzyme}{enzima} para baixar o $K_m$ dela,
ou fabricar mais \omterm{def:b1:enzymes:enzyme}{enzima}.
\textbf{20.} $K_m^{\text{ap}} = 2(1 + 0.4/0.2) = \qty{6}{mmol/L}$:
velocidade $\times 0.5/6.5$ em vez de $0.5/2.5$: \qty{38}{\%} do valor
anterior, \num{1.8e6} por segundo --- o produto estrangula a própria
síntese até abaixo da demanda, de modo que ele é consumido e o nível
dele cai, o que libera a \omterm{def:b1:enzymes:enzyme}{enzima}: um suprimento que se autorregula.
\textbf{21.} Antes: $0.5/2.5 = 0.20$ de $V_{\max}$; depois: $0.5/0.9 =
0.56$: quase o triplo.
\textbf{22.} Quase nula: em pH 5 a histidina do \omterm{def:b1:enzymes:enzyme}{sítio ativo} está
protonada e não consegue agir como base, e os resíduos ácidos do sítio
estão neutralizados; o estado de ionização da \omterm{def:b1:enzymes:enzyme}{enzima}, e talvez o
enovelamento dela, estão errados. Ela é feita para o \omterm{def:b1:eukaryotic-cell:organelle}{citosol}, e não para
o \omterm{def:b1:eukaryotic-cell:endomembrane}{lisossomo}.
\textbf{23.} O $K_m$ muda pouco (a ligação se deve em grande parte ao
resto do sítio); o $k_{\text{cat}}$ desaba por ordens de grandeza, já
que a histidina era a base catalítica que ativava a água ou a serina.
\textbf{24.} Perto do $K_m$ a velocidade responde às variações de
\omterm{def:b1:enzymes:enzyme}{substrato} (inclinação próxima de $V_{\max}/2K_m$); muito acima dele a
\omterm{def:b1:enzymes:enzyme}{enzima} fica saturada e insensível, e o excesso de \omterm{def:b1:enzymes:enzyme}{substrato} seria um
fardo osmótico e químico. Trabalhar perto do $K_m$ torna a via
controlável pela oferta.
\textbf{25.} $K_m = \qty{2.0}{mmol/L}$, $V_{\max} = \qty{1.0}{\micro mol/min}$
(\omterm{def:b1:enzymes:enzyme}{enzima} a \qty{10}{nmol/L}), $k_{\text{cat}} = \qty{1670}{s^{-1}}$; I é
competitivo, $K_i = \qty{1.0}{mmol/L}$; J é não competitivo, $K_i =
\qty{1.0}{mmol/L}$.
\end{solution}
